% --- Archon physics formalization source begin ---
% archon:physics
% archon:covers PhyXMiniProblems/problem_phyx_mini_0768.lean
% archon:source-report reports/phyx_mini/problem_phyx_mini_0768.source.json

\chapter{Physics problem phyx\_mini\_0768}
\label{ch:PhyXMiniProblems_problem_phyx_mini_0768}

\paragraph{Problem source.}
\#\# Physical scenario

In a shipping company distribution center, an open cart of mass 50.0 kg is rolling to the left at a speed of 5.00 m/s as shown in figure. Ignore friction between the cart and the floor. A 15.0 kg package slides down a chute that is inclined at 37\textbackslash{}(\textbackslash{}degree\textbackslash{}) from the horizontal and leaves the end of the chute with a speed of 3.00 m/s. The package lands in the cart and they roll together. If the lower end of the chute is a vertical distance of 4.00 m above the bottom of the cart.

\#\# Question

What are the final speed of the cart?

\#\# Answer choices

- A: A: 2.15m/s

- B: B: 3.29m/s

- C: C: 4.08m/s

- D: D: 4.14m/s

\#\# Figure caption (auxiliary; use the image as primary evidence)

The image depicts a physics scenario involving an inclined plane and a cart:

1. **Inclined Plane and Block**: A block (shown in orange) is on an inclined plane. The incline has an angle of 37° from the horizontal.

2. **Distance**: There is a vertical line with an arrow indicating a distance of 4.00 meters, suggesting the height of the block above the ground.

3. **Arrows**: Arrows show the expected motion of the block down the ramp and towards the cart, hinting at a physics problem involving gravitational pull and possibly kinetic energy transfer.

4. **Cart**: To the right of the inclined plane, there is an outline of a cart with wheels, suggesting that the block is likely meant to interact with or land in the cart following its descent.

The image is likely used to illustrate concepts such as projectile motion, energy conservation, or Newton's laws in physics problems.

\#\# Dataset metadata

Momentum and Collisions; Multi-Formula Reasoning, Physical Model Grounding Reasoning

\paragraph{Recorded answer/context.}
B: B: 3.29m/s

\paragraph{Figure/image path.}
/root/proposal\_for\_physic/science-mango/phyx\_mini\_run/phyx\_data/test\_image/768.png

\paragraph{Formalization target.}
create a compiling Lean file with sorry bodies at `PhyXMiniProblems/problem\_phyx\_mini\_0768.lean`. The Lean declarations must preserve the physical quantities, dimensions or dimensional roles, figure labels, governing-law hypotheses, and final relation expressed by this problem.
Use Mathlib/Physlib names found through LeanExplore where available. If a physics API is missing, introduce faithful local abstractions rather than scalar placeholder aliases.

\begin{theorem}[Physics formalization target]
\label{thm:physics:phyx_mini_0768:target}
\lean{PhyXMiniProblems.ProblemPhyXMini0768.final_cart_speed_after_package_capture}
\uses{def:physics:phyx-mini-0768:phyxminiproblems-problemphyxmini0768-speedinmeterspersecond, def:physics:phyx-mini-0768:phyxminiproblems-problemphyxmini0768-degreestoradians, def:physics:phyx-mini-0768:phyxminiproblems-problemphyxmini0768-shippingcartcapturesetup, def:physics:phyx-mini-0768:phyxminiproblems-problemphyxmini0768-matchesproblemreadouts, def:physics:phyx-mini-0768:phyxminiproblems-problemphyxmini0768-matchesqualitativescenario, def:physics:phyx-mini-0768:phyxminiproblems-problemphyxmini0768-matchesprimaryfigure, def:physics:phyx-mini-0768:phyxminiproblems-problemphyxmini0768-matchesphysicalgeometry, def:physics:phyx-mini-0768:phyxminiproblems-problemphyxmini0768-hasphysicalparameters, def:physics:phyx-mini-0768:phyxminiproblems-problemphyxmini0768-satisfiesinitialvelocitygeometry, def:physics:phyx-mini-0768:phyxminiproblems-problemphyxmini0768-usesstandardnearearthgravity, def:physics:phyx-mini-0768:phyxminiproblems-problemphyxmini0768-satisfiesprojectilemotion, def:physics:phyx-mini-0768:phyxminiproblems-problemphyxmini0768-satisfiesinelasticcapturelaws, def:physics:phyx-mini-0768:phyxminiproblems-problemphyxmini0768-matchesanswerchoice, def:physics:phyx-mini-0768:phyxminiproblems-problemphyxmini0768-isuniqueclosestdisplayedchoice}
The assigned autoformalize agent should translate this physics problem into Lean declarations in the covered file, with theorem and lemma proof bodies written as `by sorry`.
\end{theorem}
\begin{proof}
This is an autoformalization task, not a proof task. Produce statements that can later be proved by the physics prover without weakening the physical model.
\end{proof}

% --- Archon declaration topology begin ---
\section{Declaration topology}
\begin{definition}[velocity Dimension]
  \label{def:physics:phyx-mini-0768:phyxminiproblems-problemphyxmini0768-velocitydimension}
  \lean{PhyXMiniProblems.ProblemPhyXMini0768.velocityDimension}
  The physical dimension of velocity, L T⁻¹
\end{definition}
\begin{proof}
  This is the declared model definition of velocity Dimension.
\end{proof}
\begin{definition}[acceleration Dimension]
  \label{def:physics:phyx-mini-0768:phyxminiproblems-problemphyxmini0768-accelerationdimension}
  \lean{PhyXMiniProblems.ProblemPhyXMini0768.accelerationDimension}
  The physical dimension of acceleration, L T⁻²
\end{definition}
\begin{proof}
  This is the declared model definition of acceleration Dimension.
\end{proof}
\begin{definition}[Mass Quantity]
  \label{def:physics:phyx-mini-0768:phyxminiproblems-problemphyxmini0768-massquantity}
  \lean{PhyXMiniProblems.ProblemPhyXMini0768.MassQuantity}
  A nonnegative physical mass, independent of the chosen unit system
\end{definition}
\begin{proof}
  This is the declared model definition of Mass Quantity.
\end{proof}
\begin{definition}[Length Quantity]
  \label{def:physics:phyx-mini-0768:phyxminiproblems-problemphyxmini0768-lengthquantity}
  \lean{PhyXMiniProblems.ProblemPhyXMini0768.LengthQuantity}
  A nonnegative physical length
\end{definition}
\begin{proof}
  This is the declared model definition of Length Quantity.
\end{proof}
\begin{definition}[Time Quantity]
  \label{def:physics:phyx-mini-0768:phyxminiproblems-problemphyxmini0768-timequantity}
  \lean{PhyXMiniProblems.ProblemPhyXMini0768.TimeQuantity}
  A nonnegative physical duration
\end{definition}
\begin{proof}
  This is the declared model definition of Time Quantity.
\end{proof}
\begin{definition}[Speed Quantity]
  \label{def:physics:phyx-mini-0768:phyxminiproblems-problemphyxmini0768-speedquantity}
  \lean{PhyXMiniProblems.ProblemPhyXMini0768.SpeedQuantity}
  A nonnegative physical speed magnitude
\end{definition}
\begin{proof}
  This is the declared model definition of Speed Quantity.
\end{proof}
\begin{definition}[Signed Velocity Quantity]
  \label{def:physics:phyx-mini-0768:phyxminiproblems-problemphyxmini0768-signedvelocityquantity}
  \lean{PhyXMiniProblems.ProblemPhyXMini0768.SignedVelocityQuantity}
  \uses{def:physics:phyx-mini-0768:phyxminiproblems-problemphyxmini0768-velocitydimension}
  A signed velocity component along the horizontal floor
\end{definition}
\begin{proof}
  This is the declared model definition of Signed Velocity Quantity.
\end{proof}
\begin{definition}[Planar Position Quantity]
  \label{def:physics:phyx-mini-0768:phyxminiproblems-problemphyxmini0768-planarpositionquantity}
  \lean{PhyXMiniProblems.ProblemPhyXMini0768.PlanarPositionQuantity}
  A signed physical position vector in the plane of the figure
\end{definition}
\begin{proof}
  This is the declared model definition of Planar Position Quantity.
\end{proof}
\begin{definition}[Planar Velocity Quantity]
  \label{def:physics:phyx-mini-0768:phyxminiproblems-problemphyxmini0768-planarvelocityquantity}
  \lean{PhyXMiniProblems.ProblemPhyXMini0768.PlanarVelocityQuantity}
  \uses{def:physics:phyx-mini-0768:phyxminiproblems-problemphyxmini0768-velocitydimension}
  A signed physical velocity vector in the plane of the figure
\end{definition}
\begin{proof}
  This is the declared model definition of Planar Velocity Quantity.
\end{proof}
\begin{definition}[Planar Acceleration Quantity]
  \label{def:physics:phyx-mini-0768:phyxminiproblems-problemphyxmini0768-planaraccelerationquantity}
  \lean{PhyXMiniProblems.ProblemPhyXMini0768.PlanarAccelerationQuantity}
  \uses{def:physics:phyx-mini-0768:phyxminiproblems-problemphyxmini0768-accelerationdimension}
  A signed physical acceleration vector in the plane of the figure
\end{definition}
\begin{proof}
  This is the declared model definition of Planar Acceleration Quantity.
\end{proof}
\begin{definition}[x Axis]
  \label{def:physics:phyx-mini-0768:phyxminiproblems-problemphyxmini0768-xaxis}
  \lean{PhyXMiniProblems.ProblemPhyXMini0768.xAxis}
  Coordinate 0, horizontal and positive to the right
\end{definition}
\begin{proof}
  This is the declared model definition of x Axis.
\end{proof}
\begin{definition}[y Axis]
  \label{def:physics:phyx-mini-0768:phyxminiproblems-problemphyxmini0768-yaxis}
  \lean{PhyXMiniProblems.ProblemPhyXMini0768.yAxis}
  Coordinate 1, vertical and positive upward
\end{definition}
\begin{proof}
  This is the declared model definition of y Axis.
\end{proof}
\begin{definition}[nonnegative Readout]
  \label{def:physics:phyx-mini-0768:phyxminiproblems-problemphyxmini0768-nonnegativereadout}
  \lean{PhyXMiniProblems.ProblemPhyXMini0768.nonnegativeReadout}
  Read a nonnegative dimensionful scalar in a coherent unit system
\end{definition}
\begin{proof}
  This is the declared model definition of nonnegative Readout.
\end{proof}
\begin{definition}[signed Readout]
  \label{def:physics:phyx-mini-0768:phyxminiproblems-problemphyxmini0768-signedreadout}
  \lean{PhyXMiniProblems.ProblemPhyXMini0768.signedReadout}
  Read a signed dimensionful scalar in a coherent unit system
\end{definition}
\begin{proof}
  This is the declared model definition of signed Readout.
\end{proof}
\begin{definition}[planar Vector Readout]
  \label{def:physics:phyx-mini-0768:phyxminiproblems-problemphyxmini0768-planarvectorreadout}
  \lean{PhyXMiniProblems.ProblemPhyXMini0768.planarVectorReadout}
  Read a dimensionful planar vector in a coherent unit system
\end{definition}
\begin{proof}
  This is the declared model definition of planar Vector Readout.
\end{proof}
\begin{definition}[mass Readout]
  \label{def:physics:phyx-mini-0768:phyxminiproblems-problemphyxmini0768-massreadout}
  \lean{PhyXMiniProblems.ProblemPhyXMini0768.massReadout}
  \uses{def:physics:phyx-mini-0768:phyxminiproblems-problemphyxmini0768-massquantity, def:physics:phyx-mini-0768:phyxminiproblems-problemphyxmini0768-nonnegativereadout}
  Read a physical mass in a selected mass unit
\end{definition}
\begin{proof}
  This is the declared model definition of mass Readout.
\end{proof}
\begin{definition}[length Readout]
  \label{def:physics:phyx-mini-0768:phyxminiproblems-problemphyxmini0768-lengthreadout}
  \lean{PhyXMiniProblems.ProblemPhyXMini0768.lengthReadout}
  \uses{def:physics:phyx-mini-0768:phyxminiproblems-problemphyxmini0768-lengthquantity, def:physics:phyx-mini-0768:phyxminiproblems-problemphyxmini0768-nonnegativereadout}
  Read a physical length in a selected length unit
\end{definition}
\begin{proof}
  This is the declared model definition of length Readout.
\end{proof}
\begin{definition}[time Readout]
  \label{def:physics:phyx-mini-0768:phyxminiproblems-problemphyxmini0768-timereadout}
  \lean{PhyXMiniProblems.ProblemPhyXMini0768.timeReadout}
  \uses{def:physics:phyx-mini-0768:phyxminiproblems-problemphyxmini0768-timequantity, def:physics:phyx-mini-0768:phyxminiproblems-problemphyxmini0768-nonnegativereadout}
  Read a physical duration in a selected time unit
\end{definition}
\begin{proof}
  This is the declared model definition of time Readout.
\end{proof}
\begin{definition}[speed Readout]
  \label{def:physics:phyx-mini-0768:phyxminiproblems-problemphyxmini0768-speedreadout}
  \lean{PhyXMiniProblems.ProblemPhyXMini0768.speedReadout}
  \uses{def:physics:phyx-mini-0768:phyxminiproblems-problemphyxmini0768-speedquantity, def:physics:phyx-mini-0768:phyxminiproblems-problemphyxmini0768-nonnegativereadout}
  Read a speed in coherent selected length and time units
\end{definition}
\begin{proof}
  This is the declared model definition of speed Readout.
\end{proof}
\begin{definition}[signed Velocity Readout]
  \label{def:physics:phyx-mini-0768:phyxminiproblems-problemphyxmini0768-signedvelocityreadout}
  \lean{PhyXMiniProblems.ProblemPhyXMini0768.signedVelocityReadout}
  \uses{def:physics:phyx-mini-0768:phyxminiproblems-problemphyxmini0768-signedvelocityquantity, def:physics:phyx-mini-0768:phyxminiproblems-problemphyxmini0768-signedreadout}
  Read a signed velocity in coherent selected length and time units
\end{definition}
\begin{proof}
  This is the declared model definition of signed Velocity Readout.
\end{proof}
\begin{definition}[position Vector Readout]
  \label{def:physics:phyx-mini-0768:phyxminiproblems-problemphyxmini0768-positionvectorreadout}
  \lean{PhyXMiniProblems.ProblemPhyXMini0768.positionVectorReadout}
  \uses{def:physics:phyx-mini-0768:phyxminiproblems-problemphyxmini0768-planarpositionquantity, def:physics:phyx-mini-0768:phyxminiproblems-problemphyxmini0768-planarvectorreadout}
  Read a planar position in a selected length unit
\end{definition}
\begin{proof}
  This is the declared model definition of position Vector Readout.
\end{proof}
\begin{definition}[velocity Vector Readout]
  \label{def:physics:phyx-mini-0768:phyxminiproblems-problemphyxmini0768-velocityvectorreadout}
  \lean{PhyXMiniProblems.ProblemPhyXMini0768.velocityVectorReadout}
  \uses{def:physics:phyx-mini-0768:phyxminiproblems-problemphyxmini0768-planarvelocityquantity, def:physics:phyx-mini-0768:phyxminiproblems-problemphyxmini0768-planarvectorreadout}
  Read a planar velocity in coherent selected length and time units
\end{definition}
\begin{proof}
  This is the declared model definition of velocity Vector Readout.
\end{proof}
\begin{definition}[acceleration Vector Readout]
  \label{def:physics:phyx-mini-0768:phyxminiproblems-problemphyxmini0768-accelerationvectorreadout}
  \lean{PhyXMiniProblems.ProblemPhyXMini0768.accelerationVectorReadout}
  \uses{def:physics:phyx-mini-0768:phyxminiproblems-problemphyxmini0768-planaraccelerationquantity, def:physics:phyx-mini-0768:phyxminiproblems-problemphyxmini0768-planarvectorreadout}
  Read a planar acceleration in coherent selected length and time units
\end{definition}
\begin{proof}
  This is the declared model definition of acceleration Vector Readout.
\end{proof}
\begin{definition}[mass In Kilograms]
  \label{def:physics:phyx-mini-0768:phyxminiproblems-problemphyxmini0768-massinkilograms}
  \lean{PhyXMiniProblems.ProblemPhyXMini0768.massInKilograms}
  \uses{def:physics:phyx-mini-0768:phyxminiproblems-problemphyxmini0768-massquantity, def:physics:phyx-mini-0768:phyxminiproblems-problemphyxmini0768-massreadout}
  Kilogram readout used in the stated data and final calculation
\end{definition}
\begin{proof}
  This is the declared model definition of mass In Kilograms.
\end{proof}
\begin{definition}[length In Meters]
  \label{def:physics:phyx-mini-0768:phyxminiproblems-problemphyxmini0768-lengthinmeters}
  \lean{PhyXMiniProblems.ProblemPhyXMini0768.lengthInMeters}
  \uses{def:physics:phyx-mini-0768:phyxminiproblems-problemphyxmini0768-lengthquantity, def:physics:phyx-mini-0768:phyxminiproblems-problemphyxmini0768-lengthreadout}
  Metre readout used for the displayed vertical separation
\end{definition}
\begin{proof}
  This is the declared model definition of length In Meters.
\end{proof}
\begin{definition}[time In Seconds]
  \label{def:physics:phyx-mini-0768:phyxminiproblems-problemphyxmini0768-timeinseconds}
  \lean{PhyXMiniProblems.ProblemPhyXMini0768.timeInSeconds}
  \uses{def:physics:phyx-mini-0768:phyxminiproblems-problemphyxmini0768-timequantity, def:physics:phyx-mini-0768:phyxminiproblems-problemphyxmini0768-timereadout}
  Second readout of the package flight time
\end{definition}
\begin{proof}
  This is the declared model definition of time In Seconds.
\end{proof}
\begin{definition}[speed In Meters Per Second]
  \label{def:physics:phyx-mini-0768:phyxminiproblems-problemphyxmini0768-speedinmeterspersecond}
  \lean{PhyXMiniProblems.ProblemPhyXMini0768.speedInMetersPerSecond}
  \uses{def:physics:phyx-mini-0768:phyxminiproblems-problemphyxmini0768-speedquantity, def:physics:phyx-mini-0768:phyxminiproblems-problemphyxmini0768-speedreadout}
  Metres-per-second readout of a scalar speed
\end{definition}
\begin{proof}
  This is the declared model definition of speed In Meters Per Second.
\end{proof}
\begin{definition}[signed Velocity In Meters Per Second]
  \label{def:physics:phyx-mini-0768:phyxminiproblems-problemphyxmini0768-signedvelocityinmeterspersecond}
  \lean{PhyXMiniProblems.ProblemPhyXMini0768.signedVelocityInMetersPerSecond}
  \uses{def:physics:phyx-mini-0768:phyxminiproblems-problemphyxmini0768-signedvelocityquantity, def:physics:phyx-mini-0768:phyxminiproblems-problemphyxmini0768-signedvelocityreadout}
  Metres-per-second readout of a signed horizontal velocity
\end{definition}
\begin{proof}
  This is the declared model definition of signed Velocity In Meters Per Second.
\end{proof}
\begin{definition}[position In Meters]
  \label{def:physics:phyx-mini-0768:phyxminiproblems-problemphyxmini0768-positioninmeters}
  \lean{PhyXMiniProblems.ProblemPhyXMini0768.positionInMeters}
  \uses{def:physics:phyx-mini-0768:phyxminiproblems-problemphyxmini0768-planarpositionquantity, def:physics:phyx-mini-0768:phyxminiproblems-problemphyxmini0768-positionvectorreadout}
  Cartesian metre readout of a planar position
\end{definition}
\begin{proof}
  This is the declared model definition of position In Meters.
\end{proof}
\begin{definition}[velocity In Meters Per Second]
  \label{def:physics:phyx-mini-0768:phyxminiproblems-problemphyxmini0768-velocityinmeterspersecond}
  \lean{PhyXMiniProblems.ProblemPhyXMini0768.velocityInMetersPerSecond}
  \uses{def:physics:phyx-mini-0768:phyxminiproblems-problemphyxmini0768-planarvelocityquantity, def:physics:phyx-mini-0768:phyxminiproblems-problemphyxmini0768-velocityvectorreadout}
  Cartesian metres-per-second readout of a planar velocity
\end{definition}
\begin{proof}
  This is the declared model definition of velocity In Meters Per Second.
\end{proof}
\begin{definition}[acceleration In Meters Per Second Squared]
  \label{def:physics:phyx-mini-0768:phyxminiproblems-problemphyxmini0768-accelerationinmeterspersecondsquared}
  \lean{PhyXMiniProblems.ProblemPhyXMini0768.accelerationInMetersPerSecondSquared}
  \uses{def:physics:phyx-mini-0768:phyxminiproblems-problemphyxmini0768-planaraccelerationquantity, def:physics:phyx-mini-0768:phyxminiproblems-problemphyxmini0768-accelerationvectorreadout}
  Cartesian metres-per-second-squared readout of a planar acceleration
\end{definition}
\begin{proof}
  This is the declared model definition of acceleration In Meters Per Second Squared.
\end{proof}
\begin{definition}[degrees To Radians]
  \label{def:physics:phyx-mini-0768:phyxminiproblems-problemphyxmini0768-degreestoradians}
  \lean{PhyXMiniProblems.ProblemPhyXMini0768.degreesToRadians}
  Convert a dimensionless angle readout in degrees to radians
\end{definition}
\begin{proof}
  This is the declared model definition of degrees To Radians.
\end{proof}
\begin{definition}[Projectile Model]
  \label{def:physics:phyx-mini-0768:phyxminiproblems-problemphyxmini0768-projectilemodel}
  \lean{PhyXMiniProblems.ProblemPhyXMini0768.ProjectileModel}
  Idealization used for the airborne package
\end{definition}
\begin{proof}
  This is the declared model definition of Projectile Model.
\end{proof}
\begin{definition}[Floor Condition]
  \label{def:physics:phyx-mini-0768:phyxminiproblems-problemphyxmini0768-floorcondition}
  \lean{PhyXMiniProblems.ProblemPhyXMini0768.FloorCondition}
  Condition of the horizontal floor supporting the cart
\end{definition}
\begin{proof}
  This is the declared model definition of Floor Condition.
\end{proof}
\begin{definition}[Capture Outcome]
  \label{def:physics:phyx-mini-0768:phyxminiproblems-problemphyxmini0768-captureoutcome}
  \lean{PhyXMiniProblems.ProblemPhyXMini0768.CaptureOutcome}
  Outcome of the package--cart encounter
\end{definition}
\begin{proof}
  This is the declared model definition of Capture Outcome.
\end{proof}
\begin{definition}[Motion Direction]
  \label{def:physics:phyx-mini-0768:phyxminiproblems-problemphyxmini0768-motiondirection}
  \lean{PhyXMiniProblems.ProblemPhyXMini0768.MotionDirection}
  Qualitative directions distinguished by the two grey arrows
\end{definition}
\begin{proof}
  This is the declared model definition of Motion Direction.
\end{proof}
\begin{definition}[Figure Arrow]
  \label{def:physics:phyx-mini-0768:phyxminiproblems-problemphyxmini0768-figurearrow}
  \lean{PhyXMiniProblems.ProblemPhyXMini0768.FigureArrow}
  The two motion arrows visible in image 768.png
\end{definition}
\begin{proof}
  This is the declared model definition of Figure Arrow.
\end{proof}
\begin{definition}[Figure Element]
  \label{def:physics:phyx-mini-0768:phyxminiproblems-problemphyxmini0768-figureelement}
  \lean{PhyXMiniProblems.ProblemPhyXMini0768.FigureElement}
  Visible physical objects and geometric annotations
\end{definition}
\begin{proof}
  This is the declared model definition of Figure Element.
\end{proof}
\begin{definition}[Figure Label]
  \label{def:physics:phyx-mini-0768:phyxminiproblems-problemphyxmini0768-figurelabel}
  \lean{PhyXMiniProblems.ProblemPhyXMini0768.FigureLabel}
  Literal numerical labels printed in the supplied raster
\end{definition}
\begin{proof}
  This is the declared model definition of Figure Label.
\end{proof}
\begin{definition}[Figure Point]
  \label{def:physics:phyx-mini-0768:phyxminiproblems-problemphyxmini0768-figurepoint}
  \lean{PhyXMiniProblems.ProblemPhyXMini0768.FigurePoint}
  Distinguished endpoints used by the height annotation
\end{definition}
\begin{proof}
  This is the declared model definition of Figure Point.
\end{proof}
\begin{definition}[Shipping Cart Figure]
  \label{def:physics:phyx-mini-0768:phyxminiproblems-problemphyxmini0768-shippingcartfigure}
  \lean{PhyXMiniProblems.ProblemPhyXMini0768.ShippingCartFigure}
  \uses{def:physics:phyx-mini-0768:phyxminiproblems-problemphyxmini0768-motiondirection, def:physics:phyx-mini-0768:phyxminiproblems-problemphyxmini0768-figurearrow, def:physics:phyx-mini-0768:phyxminiproblems-problemphyxmini0768-figureelement, def:physics:phyx-mini-0768:phyxminiproblems-problemphyxmini0768-figurelabel, def:physics:phyx-mini-0768:phyxminiproblems-problemphyxmini0768-figurepoint}
  Qualitative and schematic content transcribed from the primary raster
\end{definition}
\begin{proof}
  This is the declared model definition of Shipping Cart Figure.
\end{proof}
\begin{definition}[Shipping Cart Capture Setup]
  \label{def:physics:phyx-mini-0768:phyxminiproblems-problemphyxmini0768-shippingcartcapturesetup}
  \lean{PhyXMiniProblems.ProblemPhyXMini0768.ShippingCartCaptureSetup}
  \uses{def:physics:phyx-mini-0768:phyxminiproblems-problemphyxmini0768-massquantity, def:physics:phyx-mini-0768:phyxminiproblems-problemphyxmini0768-lengthquantity, def:physics:phyx-mini-0768:phyxminiproblems-problemphyxmini0768-timequantity, def:physics:phyx-mini-0768:phyxminiproblems-problemphyxmini0768-speedquantity, def:physics:phyx-mini-0768:phyxminiproblems-problemphyxmini0768-signedvelocityquantity, def:physics:phyx-mini-0768:phyxminiproblems-problemphyxmini0768-planarpositionquantity, def:physics:phyx-mini-0768:phyxminiproblems-problemphyxmini0768-planarvelocityquantity, def:physics:phyx-mini-0768:phyxminiproblems-problemphyxmini0768-planaraccelerationquantity, def:physics:phyx-mini-0768:phyxminiproblems-problemphyxmini0768-projectilemodel, def:physics:phyx-mini-0768:phyxminiproblems-problemphyxmini0768-floorcondition, def:physics:phyx-mini-0768:phyxminiproblems-problemphyxmini0768-captureoutcome, def:physics:phyx-mini-0768:phyxminiproblems-problemphyxmini0768-shippingcartfigure}
  Independent physical quantities and intermediate states in the two-stage model. finalCartSpeed is an unconstrained observable here; no answer-choice value or solved numerical expression is stored in the setup
\end{definition}
\begin{proof}
  This is the declared model definition of Shipping Cart Capture Setup.
\end{proof}
\begin{definition}[Matches Problem Readouts]
  \label{def:physics:phyx-mini-0768:phyxminiproblems-problemphyxmini0768-matchesproblemreadouts}
  \lean{PhyXMiniProblems.ProblemPhyXMini0768.MatchesProblemReadouts}
  \uses{def:physics:phyx-mini-0768:phyxminiproblems-problemphyxmini0768-massinkilograms, def:physics:phyx-mini-0768:phyxminiproblems-problemphyxmini0768-lengthinmeters, def:physics:phyx-mini-0768:phyxminiproblems-problemphyxmini0768-speedinmeterspersecond, def:physics:phyx-mini-0768:phyxminiproblems-problemphyxmini0768-shippingcartcapturesetup}
  The five numerical givens in the prose. Neither final speed nor common final velocity is constrained here
\end{definition}
\begin{proof}
  This is the declared model definition of Matches Problem Readouts.
\end{proof}
\begin{definition}[Matches Qualitative Scenario]
  \label{def:physics:phyx-mini-0768:phyxminiproblems-problemphyxmini0768-matchesqualitativescenario}
  \lean{PhyXMiniProblems.ProblemPhyXMini0768.MatchesQualitativeScenario}
  \uses{def:physics:phyx-mini-0768:phyxminiproblems-problemphyxmini0768-shippingcartcapturesetup}
  The qualitative idealizations and encounter outcome stated in the prose
\end{definition}
\begin{proof}
  This is the declared model definition of Matches Qualitative Scenario.
\end{proof}
\begin{definition}[Matches Primary Figure]
  \label{def:physics:phyx-mini-0768:phyxminiproblems-problemphyxmini0768-matchesprimaryfigure}
  \lean{PhyXMiniProblems.ProblemPhyXMini0768.MatchesPrimaryFigure}
  \uses{def:physics:phyx-mini-0768:phyxminiproblems-problemphyxmini0768-shippingcartcapturesetup}
  Primary-image evidence: a right-descending chute and package lie above and to the left of an open cart; the package arrow points down-right, the cart arrow points left, and the 37 degrees and 4.00 m annotations are present
\end{definition}
\begin{proof}
  This is the declared model definition of Matches Primary Figure.
\end{proof}
\begin{definition}[Matches Physical Geometry]
  \label{def:physics:phyx-mini-0768:phyxminiproblems-problemphyxmini0768-matchesphysicalgeometry}
  \lean{PhyXMiniProblems.ProblemPhyXMini0768.MatchesPhysicalGeometry}
  \uses{def:physics:phyx-mini-0768:phyxminiproblems-problemphyxmini0768-xaxis, def:physics:phyx-mini-0768:phyxminiproblems-problemphyxmini0768-yaxis, def:physics:phyx-mini-0768:phyxminiproblems-problemphyxmini0768-lengthinmeters, def:physics:phyx-mini-0768:phyxminiproblems-problemphyxmini0768-positioninmeters, def:physics:phyx-mini-0768:phyxminiproblems-problemphyxmini0768-shippingcartcapturesetup}
  Physical geometry associated with the 4.00 m annotation and capture point. The package moves to the right while falling from the chute exit to the cart
\end{definition}
\begin{proof}
  This is the declared model definition of Matches Physical Geometry.
\end{proof}
\begin{definition}[Has Physical Parameters]
  \label{def:physics:phyx-mini-0768:phyxminiproblems-problemphyxmini0768-hasphysicalparameters}
  \lean{PhyXMiniProblems.ProblemPhyXMini0768.HasPhysicalParameters}
  \uses{def:physics:phyx-mini-0768:phyxminiproblems-problemphyxmini0768-xaxis, def:physics:phyx-mini-0768:phyxminiproblems-problemphyxmini0768-yaxis, def:physics:phyx-mini-0768:phyxminiproblems-problemphyxmini0768-massinkilograms, def:physics:phyx-mini-0768:phyxminiproblems-problemphyxmini0768-lengthinmeters, def:physics:phyx-mini-0768:phyxminiproblems-problemphyxmini0768-timeinseconds, def:physics:phyx-mini-0768:phyxminiproblems-problemphyxmini0768-speedinmeterspersecond, def:physics:phyx-mini-0768:phyxminiproblems-problemphyxmini0768-signedvelocityinmeterspersecond, def:physics:phyx-mini-0768:phyxminiproblems-problemphyxmini0768-velocityinmeterspersecond, def:physics:phyx-mini-0768:phyxminiproblems-problemphyxmini0768-shippingcartcapturesetup}
  Positivity and direction conditions selecting the physical branch in the figure. They contain no numerical claim about the final speed
\end{definition}
\begin{proof}
  This is the declared model definition of Has Physical Parameters.
\end{proof}
\begin{definition}[Satisfies Initial Velocity Geometry]
  \label{def:physics:phyx-mini-0768:phyxminiproblems-problemphyxmini0768-satisfiesinitialvelocitygeometry}
  \lean{PhyXMiniProblems.ProblemPhyXMini0768.SatisfiesInitialVelocityGeometry}
  \uses{def:physics:phyx-mini-0768:phyxminiproblems-problemphyxmini0768-xaxis, def:physics:phyx-mini-0768:phyxminiproblems-problemphyxmini0768-yaxis, def:physics:phyx-mini-0768:phyxminiproblems-problemphyxmini0768-speedreadout, def:physics:phyx-mini-0768:phyxminiproblems-problemphyxmini0768-signedvelocityreadout, def:physics:phyx-mini-0768:phyxminiproblems-problemphyxmini0768-velocityvectorreadout, def:physics:phyx-mini-0768:phyxminiproblems-problemphyxmini0768-degreestoradians, def:physics:phyx-mini-0768:phyxminiproblems-problemphyxmini0768-shippingcartcapturesetup}
  The stated directions and chute angle determine the initial signed cart velocity and the two components of the package's chute-exit velocity. These relations hold in every coherent choice of length and time units and contain no post-capture speed value
\end{definition}
\begin{proof}
  This is the declared model definition of Satisfies Initial Velocity Geometry.
\end{proof}
\begin{definition}[Uses Standard Near Earth Gravity]
  \label{def:physics:phyx-mini-0768:phyxminiproblems-problemphyxmini0768-usesstandardnearearthgravity}
  \lean{PhyXMiniProblems.ProblemPhyXMini0768.UsesStandardNearEarthGravity}
  \uses{def:physics:phyx-mini-0768:phyxminiproblems-problemphyxmini0768-xaxis, def:physics:phyx-mini-0768:phyxminiproblems-problemphyxmini0768-yaxis, def:physics:phyx-mini-0768:phyxminiproblems-problemphyxmini0768-accelerationinmeterspersecondsquared, def:physics:phyx-mini-0768:phyxminiproblems-problemphyxmini0768-shippingcartcapturesetup}
  The standard gravitational field used for the package's flight
\end{definition}
\begin{proof}
  This is the declared model definition of Uses Standard Near Earth Gravity.
\end{proof}
\begin{definition}[Satisfies Projectile Motion]
  \label{def:physics:phyx-mini-0768:phyxminiproblems-problemphyxmini0768-satisfiesprojectilemotion}
  \lean{PhyXMiniProblems.ProblemPhyXMini0768.SatisfiesProjectileMotion}
  \uses{def:physics:phyx-mini-0768:phyxminiproblems-problemphyxmini0768-nonnegativereadout, def:physics:phyx-mini-0768:phyxminiproblems-problemphyxmini0768-planarvectorreadout, def:physics:phyx-mini-0768:phyxminiproblems-problemphyxmini0768-shippingcartcapturesetup}
  Endpoint position and velocity equations for constant acceleration. The height datum participates through the independent physical-geometry predicate; in particular, it determines the flight time and vertical impact velocity but does not alter the conserved horizontal component
\end{definition}
\begin{proof}
  This is the declared model definition of Satisfies Projectile Motion.
\end{proof}
\begin{definition}[Satisfies Inelastic Capture Laws]
  \label{def:physics:phyx-mini-0768:phyxminiproblems-problemphyxmini0768-satisfiesinelasticcapturelaws}
  \lean{PhyXMiniProblems.ProblemPhyXMini0768.SatisfiesInelasticCaptureLaws}
  \uses{def:physics:phyx-mini-0768:phyxminiproblems-problemphyxmini0768-xaxis, def:physics:phyx-mini-0768:phyxminiproblems-problemphyxmini0768-massreadout, def:physics:phyx-mini-0768:phyxminiproblems-problemphyxmini0768-speedreadout, def:physics:phyx-mini-0768:phyxminiproblems-problemphyxmini0768-signedvelocityreadout, def:physics:phyx-mini-0768:phyxminiproblems-problemphyxmini0768-velocityvectorreadout, def:physics:phyx-mini-0768:phyxminiproblems-problemphyxmini0768-shippingcartcapturesetup}
  During the short completely inelastic capture, the frictionless floor supplies no horizontal impulse, so horizontal momentum is conserved. The floor may supply a vertical impulse; no false vertical-momentum conservation law is assumed. The last field is only the generic definition of speed as the magnitude of the common horizontal velocity
\end{definition}
\begin{proof}
  This is the declared model definition of Satisfies Inelastic Capture Laws.
\end{proof}
\begin{lemma}[package horizontal velocity at impact]
  \label{lem:physics:phyx-mini-0768:phyxminiproblems-problemphyxmini0768-package-horizontal-velocity-at-impact}
  \lean{PhyXMiniProblems.ProblemPhyXMini0768.package_horizontal_velocity_at_impact}
  \uses{def:physics:phyx-mini-0768:phyxminiproblems-problemphyxmini0768-xaxis, def:physics:phyx-mini-0768:phyxminiproblems-problemphyxmini0768-velocityinmeterspersecond, def:physics:phyx-mini-0768:phyxminiproblems-problemphyxmini0768-degreestoradians, def:physics:phyx-mini-0768:phyxminiproblems-problemphyxmini0768-shippingcartcapturesetup, def:physics:phyx-mini-0768:phyxminiproblems-problemphyxmini0768-matchesproblemreadouts, def:physics:phyx-mini-0768:phyxminiproblems-problemphyxmini0768-satisfiesinitialvelocitygeometry, def:physics:phyx-mini-0768:phyxminiproblems-problemphyxmini0768-usesstandardnearearthgravity, def:physics:phyx-mini-0768:phyxminiproblems-problemphyxmini0768-satisfiesprojectilemotion}
  Gravity has no horizontal component, so the package retains the horizontal component supplied by the chute until impact
\end{lemma}
\begin{proof}
  The conclusion follows from the governing relations and figure readouts named in the statement of package horizontal velocity at impact.
\end{proof}
\begin{lemma}[common final horizontal velocity exact]
  \label{lem:physics:phyx-mini-0768:phyxminiproblems-problemphyxmini0768-common-final-horizontal-velocity-exact}
  \lean{PhyXMiniProblems.ProblemPhyXMini0768.common_final_horizontal_velocity_exact}
  \uses{def:physics:phyx-mini-0768:phyxminiproblems-problemphyxmini0768-signedvelocityinmeterspersecond, def:physics:phyx-mini-0768:phyxminiproblems-problemphyxmini0768-degreestoradians, def:physics:phyx-mini-0768:phyxminiproblems-problemphyxmini0768-shippingcartcapturesetup, def:physics:phyx-mini-0768:phyxminiproblems-problemphyxmini0768-matchesproblemreadouts, def:physics:phyx-mini-0768:phyxminiproblems-problemphyxmini0768-satisfiesinitialvelocitygeometry, def:physics:phyx-mini-0768:phyxminiproblems-problemphyxmini0768-usesstandardnearearthgravity, def:physics:phyx-mini-0768:phyxminiproblems-problemphyxmini0768-satisfiesprojectilemotion, def:physics:phyx-mini-0768:phyxminiproblems-problemphyxmini0768-satisfiesinelasticcapturelaws}
  Substitution into horizontal momentum conservation gives the signed common velocity. The negative sign selected by the data means the loaded cart still moves left
\end{lemma}
\begin{proof}
  The conclusion follows from the governing relations and figure readouts named in the statement of common final horizontal velocity exact.
\end{proof}
\begin{definition}[Answer Choice]
  \label{def:physics:phyx-mini-0768:phyxminiproblems-problemphyxmini0768-answerchoice}
  \lean{PhyXMiniProblems.ProblemPhyXMini0768.AnswerChoice}
  Labels attached to the four final speeds printed in the problem
\end{definition}
\begin{proof}
  This is the declared model definition of Answer Choice.
\end{proof}
\begin{definition}[displayed Speed In Meters Per Second]
  \label{def:physics:phyx-mini-0768:phyxminiproblems-problemphyxmini0768-displayedspeedinmeterspersecond}
  \lean{PhyXMiniProblems.ProblemPhyXMini0768.displayedSpeedInMetersPerSecond}
  \uses{def:physics:phyx-mini-0768:phyxminiproblems-problemphyxmini0768-answerchoice}
  Displayed final speed for each choice, in metres per second
\end{definition}
\begin{proof}
  This is the declared model definition of displayed Speed In Meters Per Second.
\end{proof}
\begin{definition}[recorded Dataset Answer]
  \label{def:physics:phyx-mini-0768:phyxminiproblems-problemphyxmini0768-recordeddatasetanswer}
  \lean{PhyXMiniProblems.ProblemPhyXMini0768.recordedDatasetAnswer}
  \uses{def:physics:phyx-mini-0768:phyxminiproblems-problemphyxmini0768-answerchoice}
  The answer label recorded by the dataset; it is metadata, not a premise
\end{definition}
\begin{proof}
  This is the declared model definition of recorded Dataset Answer.
\end{proof}
\begin{definition}[Rounds To Nearest Hundredth]
  \label{def:physics:phyx-mini-0768:phyxminiproblems-problemphyxmini0768-roundstonearesthundredth}
  \lean{PhyXMiniProblems.ProblemPhyXMini0768.RoundsToNearestHundredth}
  Agreement of an actual speed with a displayed value to the nearest 0.01 m/s
\end{definition}
\begin{proof}
  This is the declared model definition of Rounds To Nearest Hundredth.
\end{proof}
\begin{definition}[Matches Answer Choice]
  \label{def:physics:phyx-mini-0768:phyxminiproblems-problemphyxmini0768-matchesanswerchoice}
  \lean{PhyXMiniProblems.ProblemPhyXMini0768.MatchesAnswerChoice}
  \uses{def:physics:phyx-mini-0768:phyxminiproblems-problemphyxmini0768-speedinmeterspersecond, def:physics:phyx-mini-0768:phyxminiproblems-problemphyxmini0768-shippingcartcapturesetup, def:physics:phyx-mini-0768:phyxminiproblems-problemphyxmini0768-answerchoice, def:physics:phyx-mini-0768:phyxminiproblems-problemphyxmini0768-displayedspeedinmeterspersecond, def:physics:phyx-mini-0768:phyxminiproblems-problemphyxmini0768-roundstonearesthundredth}
  The physical final speed rounds to the number printed by a choice
\end{definition}
\begin{proof}
  This is the declared model definition of Matches Answer Choice.
\end{proof}
\begin{definition}[Is Unique Closest Displayed Choice]
  \label{def:physics:phyx-mini-0768:phyxminiproblems-problemphyxmini0768-isuniqueclosestdisplayedchoice}
  \lean{PhyXMiniProblems.ProblemPhyXMini0768.IsUniqueClosestDisplayedChoice}
  \uses{def:physics:phyx-mini-0768:phyxminiproblems-problemphyxmini0768-speedinmeterspersecond, def:physics:phyx-mini-0768:phyxminiproblems-problemphyxmini0768-shippingcartcapturesetup, def:physics:phyx-mini-0768:phyxminiproblems-problemphyxmini0768-answerchoice, def:physics:phyx-mini-0768:phyxminiproblems-problemphyxmini0768-displayedspeedinmeterspersecond}
  The selected displayed speed is strictly closer than every other choice
\end{definition}
\begin{proof}
  This is the declared model definition of Is Unique Closest Displayed Choice.
\end{proof}
% --- Archon declaration topology end ---
% --- Archon physics formalization source end ---
